\chapter{The Unknown}

\lettrine{T}{hus} founded and recommended by its sign, the hostelry  of Master Cropole held its way steadily on towards a solid prosperity.

\zz
It was not an immense fortune that Cropole had in perspective; but he might  hope to double the thousand louis d'or left by his father, to make  another thousand louis by the sale of his house and stock, and at length to  live happily like a retired citizen.

Cropole was anxious for gain, and was half-crazy with joy at the news of the  arrival of Louis XIV.

Himself, his wife, Pittrino, and two cooks, immediately laid hands upon all the  inhabitants of the dove-cote, the poultry-yard, and the rabbit-hutches; so that  as many lamentations and cries resounded in the yards of the hostelry of the  Medici as were formerly heard in Rama.

Cropole had, at the time, but one single traveller in his house.

This was a man of scarcely thirty years of age, handsome, tall, austere, or  rather melancholy, in all his gestures and looks.

He was dressed in black velvet with jet trimmings; a white collar, as plain as  that of the severest Puritan, set off the whiteness of his youthful neck; a  small dark-coloured moustache scarcely covered his curled, disdainful lip.

He spoke to people looking them full in the face, without affectation, it is  true, but without scruple; so that the brilliancy of his black eyes became so  insupportable, that more than one look had sunk beneath his, like the weaker  sword in a single combat.

At this time, in which men, all created equal by God, were divided, thanks to  prejudices, into two distinct castes, the gentlemen and the commoner, as they  are really divided into two races, the black and the white,—at this time,  we say, he whose portrait we have just sketched could not fail of being taken  for a gentleman, and of the best class. To ascertain this, there was no  necessity to consult anything but his hands, long, slender, and white, of which  every muscle, every vein, became apparent through the skin at the least  movement, and eloquently spoke of good descent.

This gentleman, then, had arrived alone at Cropole's house. He had taken,  without hesitation, without reflection even, the principal apartment which the  hotelier had pointed out to him with a rapacious aim, very praiseworthy, some  will say, very reprehensible will say others, if they admit that Cropole was a  physiognomist, and judged people at first sight.

This apartment was that which composed the whole front of the ancient  triangular house; a large salon, lighted by two windows on the first stage, a  small chamber by the side of it, and another above it.

Now, from the time he had arrived, this gentleman had scarcely touched any  repast that had been served up to him in his chamber. He had spoken but two  words to the host, to warn him that a traveller of the name of Parry would  arrive, and to desire that, when he did, he should be shown up to him  immediately.

He afterwards preserved so profound a silence, that Cropole was almost  offended, so much did he prefer people who were good company.

This gentleman had risen early the morning of the day on which this history  begins, and had placed himself at the window of his salon, seated upon the  ledge, and leaning upon the rail of the balcony, gazing sadly but persistently  on both sides of the street, watching, no doubt, for the arrival of the  traveller he had mentioned to the host.

In this way he had seen the little cortège of Monsieur return from hunting,  then had again partaken of the profound tranquillity of the street, absorbed in  his own expectations.

All at once the movement of the crowd going to the meadows, couriers setting  out, washers of pavement, purveyors of the royal household, gabbling,  scampering shop-boys, chariots in motion, hair-dressers on the run, and pages  toiling along, this tumult and bustle had surprised him, but without losing any  of that impassible and supreme majesty which gives to the eagle and the lion  that serene and contemptuous glance amidst the hurrahs and shouts of hunters or  the curious.

Soon the cries of the victims slaughtered in the poultry-yard, the hasty steps  of Madame Cropole up that little wooden staircase, so narrow and so echoing;  the bounding pace of Pittrino, who only that morning was smoking at the door  with all the phlegm of a Dutchman; all this communicated something like  surprise and agitation to the traveller.

As he was rising to make inquiries, the door of his chamber opened. The unknown  concluded they were about to introduce the impatiently expected traveller, and  made three precipitate steps to meet him.

But, instead of the person he expected, it was Master Cropole who appeared, and  behind him, in the half-dark staircase, the pleasant face of Madame Cropole,  rendered trivial by curiosity. She only gave one furtive glance at the handsome  gentleman, and disappeared.

Cropole advanced, cap in hand, rather bent than bowing.

A gesture of the unknown interrogated him, without a word being pronounced.

<Monsieur,> said Cropole, <I come to ask how—what ought  I to say: your lordship, monsieur le comte, or monsieur le marquis?>

<Say monsieur, and speak quickly,> replied the unknown, with that  haughty accent which admits of neither discussion nor reply.

<I came, then, to inquire how monsieur had passed the night, and if  monsieur intended to keep this apartment?>

<Yes.>

<Monsieur, something has happened upon which we could not reckon.>

<What?>

<His majesty Louis XIV will enter our city to-day, and will remain here  one day, perhaps two.>

Great astonishment was painted on the countenance of the unknown.

<The King of France is coming to Blois?>

<He is on the road, monsieur.>

<Then there is the stronger reason for my remaining,> said the  unknown.

<Very well; but will monsieur keep all the apartments?>

<I do not understand you. Why should I require less to-day than  yesterday?>

<Because, monsieur, your lordship will permit me to say, yesterday I did  not think proper, when you chose your lodging, to fix any price that might have  made your lordship believe that I prejudged your resources; whilst  to-day\longdash>

The unknown coloured; the idea at once struck him that he was supposed to be  poor, and was being insulted.

<Whilst to-day,> replied he, coldly, <you do not  prejudge.>

<Monsieur, I am a well-meaning man, thank God! and simple hotelier as I  am, there is in me the blood of a gentleman. My father was a servant and  officer of the late Marechal d'Ancre. God rest his soul!>

<I do not contest that point with you; I only wish to know, and that  quickly, to what your questions tend?>

<You are too reasonable, monsieur, not to comprehend that our city is  small, that the court is about to invade it, that the houses will be  overflowing with inhabitants, and that lodgings will consequently obtain  considerable prices.>

Again the unknown coloured. <Name your terms,> said he.

<I name them with scruple, monsieur, because I seek an honest gain, and  that I wish to carry on my business without being uncivil or extravagant in my  demands. Now the room you occupy is considerable, and you are alone.>

<That is my business.>

<Oh! certainly. I do not mean to turn monsieur out.>

The blood rushed to the temples of the unknown; he darted at poor Cropole, the  descendant of one of the officers of the Marechal d'Ancre, a glance that  would have crushed him down to beneath that famous chimney-slab, if Cropole had  not been nailed to the spot by the question of his own proper interests.

<Do you desire me to go?> said he. <Explain  yourself—but quickly.>

<Monsieur, monsieur, you do not understand me. It is very  critical—I know—that which I am doing. I express myself badly, or  perhaps, as monsieur is a foreigner, which I perceive by his  accent\longdash>

In fact, the unknown spoke with that impetuosity which is the principal  character of English accentuation, even among men who speak the French language  with the greatest purity.

<As monsieur is a foreigner, I say, it is perhaps he who does not catch  my exact meaning. I wish for monsieur to give up one or two of the apartments  he occupies, which would diminish his expenses and ease my conscience. Indeed,  it is hard to increase unreasonably the price of the chambers, when one has had  the honour to let them at a reasonable price.>

<How much does the hire amount to since yesterday?>

<Monsieur, to one louis, with refreshments and the charge for the  horse.>

<Very well; and that of to-day?>

<Ah! there is the difficulty. This is the day of the king's  arrival; if the court comes to sleep here, the charge of the day is reckoned.  From that it results that three chambers, at two louis each, make six louis.  Two louis, monsieur, are not much; but six louis make a great deal.>

The unknown, from red, as we have seen him, became very pale.

He drew from his pocket, with heroic bravery, a purse embroidered with a  coat-of-arms, which he carefully concealed in the hollow of his hand. This  purse was of a thinness, a flabbiness, a hollowness, which did not escape the  eye of Cropole.

The unknown emptied the purse into his hand. It contained three double louis,  which amounted to the six louis demanded by the host.

But it was seven that Cropole had required.

He looked, therefore, at the unknown, as much as to say, <And  then?>

<There remains one louis, does there not, master hotelier?>

<Yes, monsieur, but\longdash>

The unknown plunged his hand into the pocket of his haut-de-chausses, and  emptied it. It contained a small pocket-book, a gold key, and some silver. With  this change, he made up a louis.

<Thank you, monsieur,> said Cropole. <It now only remains for  me to ask whether monsieur intends to occupy his apartments to-morrow, in which  case I will reserve them for him; whereas, if monsieur does not mean to do so,  I will promise them to some of the king's people who are coming.>

<That is but right,> said the unknown, after a long silence;  <but as I have no more money, as you have seen, and as I yet must retain  the apartments, you must either sell this diamond in the city, or hold it in  pledge.>

Cropole looked at the diamond so long, that the unknown said, hastily:

<I prefer your selling it, monsieur; for it is worth three hundred  pistoles. A Jew—are there any Jews in Blois?—would give you two  hundred or a hundred and fifty for it—take whatever may be offered for  it, if it be no more than the price of your lodging. Begone!>

<Oh! monsieur,> replied Cropole, ashamed of the sudden inferiority  which the unknown reflected upon him by this noble and disinterested  confidence, as well as by the unalterable patience opposed to so many  suspicions and evasions. <Oh, monsieur, I hope people are not so  dishonest at Blois as you seem to think; and that the diamond, being worth what  you say\longdash>

The unknown here again darted at Cropole one of his withering glances.

<I really do not understand diamonds, monsieur, I assure you,>  cried he.

<But the jewelers do: ask them,> said the unknown. <Now I  believe our accounts are settled, are they not, monsieur l'hote?>

<Yes, monsieur, and to my profound regret; for I fear I have offended  monsieur.>

<Not at all!> replied the unknown, with ineffable majesty.

<Or have appeared to be extortionate with a noble traveller. Consider,  monsieur, the peculiarity of the case.>

<Say no more about it, I desire; and leave me to myself.>

Cropole bowed profoundly, and left the room with a stupefied air, which  announced that he had a good heart, and felt genuine remorse.

The unknown himself shut the door after him, and, when left alone, looked  mournfully at the bottom of the purse, from which he had taken a small silken  bag containing the diamond, his last resource.

He dwelt likewise upon the emptiness of his pockets, turned over the papers in  his pocket-book, and convinced himself of the state of absolute destitution in  which he was about to be plunged.

He raised his eyes towards heaven, with a sublime emotion of despairing  calmness, brushed off with his hand some drops of sweat which trickled over his  noble brow, and then cast down upon the earth a look which just before had been  impressed with almost divine majesty.

That the storm had passed far from him, perhaps he had prayed in the bottom of  his soul.

He drew near to the window, resumed his place in the balcony, and remained  there, motionless, annihilated, dead, till the moment when, the heavens  beginning to darken, the first flambeaux traversed the enlivened street, and  gave the signal for illumination to all the windows of the city.  